% !TEX program = xelatex
\documentclass[UTF8, a4paper, 12pt, fontset=none]{ctexart}

% === import ===
\usepackage[margin=2.5cm]{geometry}
\usepackage{amsmath, amssymb}
\usepackage{enumitem}
\usepackage{setspace}
\usepackage{xcolor}

% === style ===
\setCJKmainfont{Noto Serif CJK TC} % 字體
\onehalfspacing % 1.5 倍行距
\pagestyle{plain} % 取消左上文字

% === title ===
\title{資料庫系統期末專題 Report 1\\ \textbf{錢錢錢市－股票模擬系統}}
\author{
    \begin{tabular}{c}
        第八組 \\ 
        41347001S 楊凱臣 \quad 41347009S 周聖詠 \\ 
        41347011S 許育綸 \quad 41347041S 盧人豪
    \end{tabular}
}
\date{}

% === content ===
\begin{document}

    \maketitle

    % === section === 1
    \section{專題動機與目標}
    在這個日新月異的時代，愈來愈多人想學習股票交易，但又不知道該如何下手。
    雖然市面上已有部分模擬交易平台，但多數僅提供即時或近期資料，
    缺乏讓使用者回到特定歷史時間點、以真實市場情境反覆驗證判斷的環境。

    因此本專題希望設計一套以台灣股市真實歷史資料為基礎的股票模擬系統，
    讓使用者選擇過去的時間點，透過實際下單與觀察 K 線、均線等技術指標，
    在零風險的環境下反覆驗證自己對市場走勢的判斷，建立對股票交易的直覺與理解。

    本系統的目標如下：
    \begin{enumerate}
        \item 提供以真實歷史資料為基礎的虛擬交易環境
        \item 讓使用者體驗完整的交易流程，包含下單、撮合與損益計算
        \item 透過反覆模擬，協助使用者建立對 K 線與市場走勢的判斷能力
    \end{enumerate}

    % === section === 2
    \section{系統使用者}
    \begin{enumerate}
        \item \textbf{一般使用者}：對股票交易有興趣的新手或初學者。
            可註冊帳號、建立模擬存檔、進行模擬、
            進行虛擬股票交易，並查詢持股與損益紀錄。
    \end{enumerate}

    % 如果要有不同使用者，我有兩個想法：
    % \section{系統使用者}
    %     本系統的使用者分為以下兩類：
    %     \begin{enumerate}
    %         \item \textbf{一般使用者}：對股票交易有興趣的新手或初學者。
    %         可註冊帳號、建立模擬存檔、選擇近五年歷史時間點開始模擬、
    %         進行虛擬股票交易，並查詢持股與損益紀錄。
    %         自選股上限為 20 支，存檔數量上限為 3 個。

    %         \item \textbf{VIP 使用者}：訂閱進階方案的使用者。
    %         擁有一般使用者的所有功能，另可選擇更完整的歷史時間範圍（2018 年至今），
    %         自選股上限提升為 100 支，存檔數量提升為 10 個。
    %     \end{enumerate}

    % \section{系統使用者}
    % 本系統的使用者分為以下兩類：
    % \begin{enumerate}
    %     \item \textbf{一般使用者}：對股票交易有興趣的新手或初學者。
    %     可註冊帳號、建立模擬存檔、選擇歷史時間點開始模擬、
    %     進行虛擬股票交易，並查詢持股與損益紀錄。

    %     \item \textbf{管理員}：負責維護系統資料與帳號管理。
    %     可查看所有使用者帳號狀態與模擬紀錄，並管理帳號權限。
    %     （管理員功能列為擴展項目，初期開發以一般使用者功能為主。）
    % \end{enumerate}


    % === section === 3
    % \section{系統主要功能}
    % \begin{enumerate}
    %     \item \textbf{使用者資料管理}：系統能創建並記錄使用者的帳號、密碼、庫存與餘額。作為股票交易的基礎。
    %     \item \textbf{股票資訊查詢}：系統能記錄每一時間段各股票的代碼、名稱與當前股價，供使用者查詢。
    %     \item \textbf{模擬交易}：使用者可進行模擬買賣股票。
    %     \item \textbf{時間管理系統}：系統可指派使用者一時間段，即時更新使用者當前時間段，並記錄時間流逝。
    % \end{enumerate}

    \section{系統主要功能}
    \begin{enumerate}
        \item \textbf{帳戶管理}：使用者可註冊與登入帳號，
        系統記錄使用者身份資訊作為交易的基礎。

        \item \textbf{模擬存檔管理}：使用者可建立多個獨立的模擬存檔，
        每個存檔擁有各自的時間軸、資金與持股狀態，互不影響。

        \item \textbf{時間軸控制}：
        系統會隨機選擇一個過去的歷史日期作為模擬起點，
        使用者能手動推進模擬時間，每次推進一個交易日，
        依序經歷盤前、開盤、盤中、收盤、盤後五個階段。

        \item \textbf{資金管理}：每個存檔包含儲蓄戶與交割戶，
        使用者可在兩者之間進行轉帳，股票交易僅從交割戶扣款。

        \item \textbf{股票查詢}：使用者可透過搜尋或類股分類瀏覽股票，
        並查看歷史 K 線圖與均線等技術指標。

        \item \textbf{自選股管理}：使用者可將感興趣的股票加入自選清單，
        方便快速追蹤。

        \item \textbf{下單與撮合}：使用者可提交限價單或市價單進行買賣，
        系統依照當日開盤價或收盤 OHLC 進行撮合。
        未成交的委託可查詢或撤銷。

        \item \textbf{持股與損益查詢}：使用者可查看目前持有的股票明細，
        包含庫存量、平均成本與未實現損益，
        以及歷史交易紀錄與已實現損益。

        \item \textbf{模擬回顧}：模擬結束後，使用者可查看整段模擬期間的
        資產走勢、交易紀錄與損益總結，了解自己的操作成效。
    \end{enumerate}


    % % === section === 4
    % \section{需求功能分析}
    % \textbf{4.1 功能性需求}
    % \begin{enumerate}[leftmargin=2cm]
    %     \item {能建立與維護使用者資料。}
    %     \item {能提供使用者查看個人庫存與餘額。}
    %     \item {能查看股票代碼、名稱與當前股價。}
    %     \item {使用者可進行模擬買賣股票。}
    %     \item {能管理系統之虛擬時間。}
    % \end{enumerate}

    % \textbf{4.2 功能性需求}
    % \begin{enumerate}[leftmargin=2cm]
    %     \item {保持資料正確性，避免重複使用者名稱。}
    %     \item {保持資料安全性，避免使用弱密碼。}
    %     \item {讓使用者能清楚查詢股票及個人之各項資訊。}
    % \end{enumerate}
    
    % === section === 4
    \section{需求功能分析}
    \subsection{功能性需求}
    \begin{enumerate}
        \item 使用者可註冊與登入帳號。
        \item 使用者可建立與管理多個模擬存檔。
        \item 系統隨機選擇歷史時間點開始模擬
        \item 使用者可手動推進模擬交易日。
        \item 使用者可在儲蓄戶與交割戶之間進行轉帳。
        \item 使用者可搜尋股票、瀏覽類股分類與查看 K 線圖。
        \item 使用者可管理自選股清單。
        \item 使用者可提交限價單或市價單，並查詢與撤銷待成交委託。
        \item 系統依照當日真實歷史價格自動撮合委託單。
        \item 使用者可查詢持股明細與損益紀錄。
        \item 使用者可在模擬結束後查看整體模擬回顧。
    \end{enumerate}

    \subsection{非功能性需求}
    \begin{enumerate}
        \item \textbf{Security}：使用者名稱須唯一，系統須拒絕重複註冊。
        \item \textbf{Security}：密碼須符合最低強度要求。
        \item \textbf{Security}：使用者只能存取自己的模擬存檔與資料。
        \item \textbf{Reliability}：交易撮合須以 Transaction 處理，確保資料一致性。
        \item \textbf{Reliability}：交割戶餘額不足時系統須拒絕成交，餘額不得為負數。
        \item \textbf{Reliability}：同一委託單不得被撮合兩次。
        \item \textbf{Maintainability}：各模擬存檔之間的資料須完全獨立，互不影響。
    \end{enumerate}


    % % === section === 5
    % \section{資料分析}
    % \begin{enumerate}
    %     \item \textbf{使用者}：用來記錄使用者資訊，包括 帳號ID、帳號、密碼、餘額、股票庫存等。
    %     \item \textbf{股票}：用來記錄股票資訊，包括 代碼、名稱、股價、時間、一系列股票指數等。
    %     \item \textbf{交易事件}：用來記錄每筆交易的內容，包括 交易 ID、時間、買賣、金額、帳號ID、股票代碼。
    %     \item \textbf{金錢記錄}：用來記錄不同時機點用戶的金錢，包括 帳號ID、時間、金錢。
    % \end{enumerate}

    % === section === 5
    \section{資料分析}
    本系統初步分析後，主要資料對象如下：

    \begin{enumerate}
        \item \textbf{使用者}：記錄帳號資訊與身份，
        作為登入驗證的基礎。

        \item \textbf{模擬存檔}：記錄一組模擬的狀態，
        包含目前所在的模擬日期與起始設定。
        每位使用者可擁有多個獨立的存檔。

        \item \textbf{帳戶}：每個存檔擁有儲蓄戶與交割戶，
        分別記錄各自的餘額。

        \item \textbf{股票}：記錄台股上市股票的基本資料，
        如代號、名稱、所屬類股。

        \item \textbf{歷史價格}：記錄每支股票每個交易日的
        開高低收價格與成交量。

        \item \textbf{委託單}：記錄使用者在某個模擬日提交的
        買賣委託，包含委託類型與當前狀態。

        \item \textbf{交易紀錄}：記錄委託單成交後的結果，
        包含成交價、成交量與相關費用。
        此資料描述委託單成交這件事本身。

        \item \textbf{持股}：記錄某個存檔目前持有某支股票的
        數量與平均成本。
        此資料描述存檔與股票之間的關聯本身。

        \item \textbf{自選股}：記錄某個存檔所關注的股票清單。
        此資料描述存檔與股票之間的關注關聯。
        
        % === 股利，看要不要做 ===
        % \item \textbf{除權息資料}：記錄每支股票歷年的現金股利資訊，
        % 供系統在模擬推進至除息日時自動計算並發放股利。
    \end{enumerate}

    % % === section === 6
    % \section{資料關聯分析}
    % \begin{enumerate}
    %     \item 一位使用者可買多支股票。
    %     \item 一支股票可有多位使用者購買。
    %     \item 使用者與股票之間存在多對多關係，這個關係可由交易事件表示。
    %     \item 交易事件需記錄交易時間、交易類別、與交易金額等資訊。
    %     \item 若加入金額記錄，可快速回朔與觀察不同時間點的金額變動。
    % \end{enumerate}
    % === section === 6
    \section{資料關聯分析}
    \begin{enumerate}
        \item 一位使用者可擁有多個模擬存檔，
        每個存檔的資料互相獨立。

        \item 每個模擬存檔擁有兩個帳戶（儲蓄戶與交割戶），
        帳戶餘額隨轉帳與交易而變動。

        \item 一個存檔可持有多支股票，一支股票也可被多個存檔持有，
        兩者之間存在多對多關係，透過持股來表示。
        持股除了連結存檔與股票之外，還需記錄持有數量與平均成本。

        \item 一個存檔可關注多支股票，一支股票也可被多個存檔關注，
        兩者之間存在多對多關係，透過自選股來表示。

        \item 一個存檔可有多筆委託單，委託單成交後產生對應的交易紀錄。
        交易紀錄除了連結委託單之外，還需記錄成交價、成交量與費用。

        \item 每支股票對應多筆歷史價格，每個交易日一筆。

        % === 股利，看要不要做 ===
        % \item 每支股票可有多筆除權息資料，每筆記錄該年度的現金股利資訊。
        % 系統推進至除息日時，自動計算並發放股利至對應存檔的帳戶。
    \end{enumerate}

    \section{運作規則}
    \begin{enumerate}
        \item 使用者須手動推進模擬時間，系統不會自動前進，
        每次推進一個交易日。

        \item 委託單僅支援當日有效(ROD)，
        收盤後未成交的委託單自動失效。

        \item 股票買賣最小單位為 1,000 股（一張），不接受零股。

        \item 買進委託僅從交割戶扣款，
        若交割戶餘額不足則無法成交。

        \item 限價買單：買入價格須高於或等於當日最低價才可成交。\\
            限價賣單：賣出價格須低於或等於當日最高價才可成交。\\
            市價單：統一以收盤價成交。
        
        \item 使用者提交委託單時，委託價格須在當日漲跌停範圍內。
        系統依據前一交易日的收盤資料取得當日漲跌停價格，
        超出範圍的委託單將被拒絕。
        
        \item 手續費為買賣各 0.1425\%（未滿 20 元以 20 元計），
        證交稅於賣出時收取（股票 0.3\%、ETF 0.1\%）。

        % === 股利，看要不要做 ===
        % \item 系統推進至除息日時，自動計算並發放現金股利至對應存檔的儲蓄戶。

        \item 每個模擬存檔的時間軸與資產狀態互相獨立，互不影響。
    \end{enumerate}
\end{document}